
\section{早期版本中的若干错误及其更正}
\label{sec:a1-6}

\noindent\textit{A1.6.i).}
在定义~\ref{def:a1-1} 中，我们只把平凡丛
$E_p\times V_p\to V_p$ 作为障碍丛来考虑；作为等变丛，它未必平凡。
在 \textbf{[FuOn99II]} 中，我们允许更一般的丛。然而，这两个定义是
等价的，因为任何丛都局部平凡，而一点的仓西邻域可以取得任意小。

\medskip
\noindent\textit{A1.6.ii).}
在本书的预印本版本 \textbf{[FOOO00]} 中，我们没有使用“轨形的光滑嵌入”
一语来重述 \textbf{[FuOn99II]} 的定义 5.3，因而定义~\ref{def:a1-3}
表述得不够严格。在条件 (A1.4.2) 中，应要求所诱导的映射
\[
 V_{pq}/\Gamma_q\longrightarrow V_p/\Gamma_p
\]
为单射；我们还加入了条件 (A1.4.6)。注意，若
$h_{pq}$-等变光滑嵌入
$\phi_{pq}\colon V_{pq}\to V_p$ 按 \textbf{[FuOn99II]} 第 941 页的意义
诱导轨形的光滑嵌入，则条件 (A1.4.6) 成立。因此，上述定义
\ref{def:a1-3}（包括条件 (A1.4.6)）同 \textbf{[FuOn99II]} 的定义 5.3
等价。

\medskip
\noindent\textit{A1.6.iii).}
在定义~\ref{def:a1-5} 的条件 (A1.6.2) 中，公式含有
$\gamma_{pqr}$。\textbf{[FOOO00]} 的式 (A2.1.11) 漏掉了这一项；
例~\ref{ex:a1-65} 说明这是一个错误。另一方面，对应的
\textbf{[FuOn99II]} 定义 5.3 本身是正确的。

\medskip
\noindent\textit{A1.6.iv).}
我们指出，定义~\ref{def:a1-14} 与 \textbf{[FuOn99II]} 中的对应定义
略有不同。后者不要求同构~\eqref{eq:a1-15} 由仓西映射的微分诱导，
只要求存在使图式~\eqref{diag:a1-16} 交换的这种同构。事实上，
\textbf{[FuOn99II]} 中定理 6.4 的证明已经隐含地使用了如下条件：
同构~\eqref{eq:a1-15} 由 $d_{\mathrm{fiber}}s_p$ 诱导。

因此，旧定义确有错误，而定义~\ref{def:a1-14} 才是切丛的正确定义
（见第~\ref{sec:a1-5} 小节）。例~\ref{ex:a1-64} 给出的空间 $X$ 带有
仓西结构；按本书的意义，它没有切丛，而按 \textbf{[FuOn99II]} 的意义
却有切丛。为了使虚基本闭链的配边不变性成立，必须排除例
\ref{ex:a1-64} 中的 $X$。

在 \textbf{[FuOn99II]} 与本书把有切丛仓西结构应用于模空间问题的所有
情形中，同构~\eqref{eq:a1-15} 都由仓西映射的微分诱导。因此，这个错误
并未造成任何问题。

在相容多值截面存在性的证明（定理~\ref{thm:a1-23}）中，我们以相容方式
延拓多值截面，也就是要求定义~\ref{def:a1-21} 中的相容条件。原仓西映射
只有在同构~\eqref{eq:a1-15} 由仓西映射诱导时，才按该定义相容。若原仓西
映射不相容，就不可能把它扰动为相容映射。因此，为了证明定理
\ref{thm:a1-23}，同构~\eqref{eq:a1-15} 必须由仓西映射诱导。
\textbf{[FuOn99II]} 的式 (6.4.4) 也要求同一相容条件。

再看第~\ref{sec:7-1} 节构造的仓西结构：法丛与障碍丛的商，都同式
\eqref{eq:7-1-16} 中出现的两种 $E$ 之选择的差同构。因此，同构
\eqref{eq:a1-15} 由 $s$ 的微分给出。
\translatornote{源书把本段首次出现的文献号排作 \textbf{[FuOnII99]}，
而全书书目键及同页其余三处均为 \textbf{[FuOn99II]}，译文据此统一；
又在后三处把同构 (A1.15) 误写成 (A1.5)，译文依本段持续讨论的切丛
同构、图式 (A1.16) 及定义 A1.21 恢复。见疑误 A-0715、A-0716。}
